% Part: normal-modal-logic
% Chapter: frame-correspondence
% Section: equivalence-S5

\documentclass[../../../include/open-logic-section]{subfiles}

\begin{document}

\olfileid{nml}{frd}{es5}

\olsection{तुल्यता संबंध आणि \Log{S5}}

सर्व वैश्विक चौकटींमध्ये वैध असलेल्या !!{formula}s
चा संच म्हणून मोडल तर्कशास्त्र~\Log{S5} निश्चित होते;
अशा चौकटीत प्रत्येक जग स्वतःसह प्रत्येक जगापासून
प्राप्य असते. येथे $\Box$ आवश्यकतेशी आणि
$\Diamond$ शक्यतेशी जुळतो: $!A$ \emph{प्रत्येक}
जगात सत्य असेल तर $\Box !A$ सत्य असते, आणि
$!A$ \emph{किमान एका} जगात सत्य असेल तर
$\Diamond !A$ सत्य असते. \Log{S5} हे सर्व परावर्ती,
सममित आणि संक्रमक चौकटींमध्ये, म्हणजे सर्व
\emph{तुल्यता संबंधांवर}, वैध असलेल्या !!{formula}s
चा संच म्हणूनही सांगता येते.

\begin{defn}
  $W$ वरील द्विस्थानी संबंध $R$ हा \emph{तुल्यता संबंध}
  असतो तेव्हा आणि केवळ तेव्हाच, जेव्हा तो परावर्ती,
  सममित आणि संक्रमक असतो. $W$ वरील संबंध $R$
  \emph{वैश्विक} असतो तेव्हा आणि केवळ तेव्हाच, जेव्हा
  सर्व $u,v \in
  W$ साठी $Ruv$ असते.
\end{defn}

\Ax{T}, \Ax{B}, आणि \Ax{4} हे अनुक्रमे परावर्ती,
सममित आणि संक्रमक चौकटींचे वैशिष्ट्य दर्शवतात.
म्हणून प्राप्यता संबंध तुल्यता संबंध असलेल्या चौकटी
नेमक्या त्या आहेत ज्यांत ही तिन्ही !!{formula}s वैध
आहेत. $R$ चे इतर परिचित गुणधर्म एकत्र केल्यास
तुल्यता संबंध परिभाषित करणाऱ्या अटींच्या समतुल्य
अटी मिळतात; म्हणून !!{formula}s च्या इतर संयोगांनीही
तुल्यता संबंधांचे वैशिष्ट्य सांगता येते.

\begin{prop}\ollabel{prop:equivalences}
  पुढील अटी समतुल्य आहेत:
  \begin{enumerate}
  \item $R$ हा तुल्यता संबंध आहे;
  \item $R$ परावर्ती आणि युक्लिडी आहे;
  \item $R$ सर्वत्र उत्तरवर्ती, सममित आणि युक्लिडी आहे;
  \item $R$ सर्वत्र उत्तरवर्ती, सममित आणि संक्रमक आहे.
  \end{enumerate}
\end{prop}

\begin{proof}
  सरावासाठी.
\end{proof}

\begin{prob}
  पुढील विधाने दाखवून \olref[nml][frd][es5]{prop:equivalences}
  सिद्ध करा:
  \begin{enumerate}
  \item $R$ सममित आणि संक्रमक असेल, तर तो युक्लिडी आहे.
  \item $R$ परावर्ती असेल, तर तो सर्वत्र उत्तरवर्ती आहे.
  \item $R$ परावर्ती आणि युक्लिडी असेल, तर तो सममित आहे.
  \item $R$ सममित आणि युक्लिडी असेल, तर तो संक्रमक आहे.
  \item $R$ सर्वत्र उत्तरवर्ती, सममित आणि संक्रमक असेल,
    तर तो परावर्ती आहे.
  \end{enumerate}
  या विधानांवरून सर्व अटी समतुल्य आहेत, हे कसे
  सिद्ध होते ते स्पष्ट करा.
\end{prob}

\olref{prop:equivalences} हे \olref[prf][prs]{prop:S5}
चे चिन्हार्थविषयक प्रतिरूप आहे: $R$ तुल्यता संबंध
असलेल्या चौकटींच्या मोडल तर्कशास्त्राचे, परंपरेने
\Log{S5} म्हणतात त्याचे, ते समतुल्य वैशिष्ट्यदर्शन देते.

वैश्विक संबंध आणि तुल्यता संबंध यांचे नाते काय?
प्रत्येक वैश्विक संबंध तुल्यता संबंध असला, तरी
प्रत्येक तुल्यता संबंध वैश्विक असेलच असे नाही.
तरी सर्व वैश्विक संबंधांवर वैध असलेली !!{formula}s
आणि सर्व तुल्यता संबंधांवर वैध असलेली सूत्रे
नेमकी तीच आहेत.

\begin{prop}
  $R$ हा तुल्यता संबंध असू द्या. प्रत्येक $w \in W$
  साठी $w$ चा \emph{तुल्यतावर्ग} म्हणजे
  $[w] = \{w'\in W :
  Rww'\}$ हा संच ठरवा. मग:
  \begin{enumerate}
  \item $w \in [w]$;
  \item प्रत्येक तुल्यतावर्ग~$[w]$ वर $R$ वैश्विक आहे;
  \item तुल्यतावर्ग $W$ ला परस्पर विसंयुक्त आणि
    एकत्रितपणे संपूर्ण W व्यापणाऱ्या उपसंचांत विभागतात.
  \end{enumerate}
\end{prop}

\begin{prop}\ollabel{prop:S5=univ}
  !!^a{formula}~$!A$ हे $R$ तुल्यता संबंध असलेल्या
  सर्व चौकटींमध्ये $\mModel{F} =\tuple{W,R}$ वैध
  असते तेव्हा आणि केवळ तेव्हाच, जेव्हा ते $R$
  वैश्विक असलेल्या सर्व चौकटींमध्ये
  $\mModel{F} =\tuple{W,R}$ वैध असते. म्हणून
  वैश्विक चौकटींचे तर्कशास्त्र म्हणजे \Log{S5}.
\end{prop}

\begin{proof}
  $W$ वरील वैश्विक संबंध~$R$ हा तुल्यता संबंध असतो,
  हे लगेच तपासता येते. म्हणून $!A$ हे $R$ तुल्यता
  संबंध असलेल्या सर्व चौकटींमध्ये वैध असेल, तर ते
  सर्व वैश्विक चौकटींमध्येही वैध असते. उलट दिशेसाठी
  व्यत्यासाभासाने युक्तिवाद करू. $W$ वरील $R$
  तुल्यता संबंध असलेल्या चौकट~$\tuple{W,R}$ वर
  आधारित $\mModel{M} = \tuple{W, R, V}$ या
  प्रतिमानात $!B$ हे !!a{formula} एखाद्या जगात~$w$
  असत्य आहे असे समजा. म्हणजे $\mSat/{M}{!B}[w]$.
  पुढीलप्रमाणे $\mModel{M}' = \tuple{W',
    R', V'}$ हे प्रतिमान ठरवा:
  \begin{enumerate}
  \item $W' = [w]$;
  \item $R'$ हा $W'$ वर वैश्विक आहे;
  \item $V'(p) = V(p) \cap W'$.
  \end{enumerate}
  (\olref{fig:partition} मधील छायांकित भाग
  $\mModel{M}'$ चा जगांचा संच~$W'$ दाखवतो.)
  $W'$ वर $R$ आणि $R'$ हे संबंध जुळतात, हे सहज
  दिसते. मग !!{formula}s वर विगमन करून, प्रत्येक
  $!A$ आणि सर्व $w' \in W'$ साठी
  $\mSat{M'}{!A}[w']$ तेव्हा आणि केवळ तेव्हाच,
  जेव्हा $\mSat{M}{!A}[w']$, हे दाखवता येते
  (कारण $W'
  \subseteq W$). विशेषतः $\mSat/{M'}{!B}[w]$;
  म्हणून $!B$ हे वैश्विक चौकटीवर आधारित
  प्रतिमानात असत्य आहे.
\end{proof}

\begin{figure}[t]
  \centering
  \begin{tikzpicture}[node distance=2cm, auto, thick]
    \clip [rounded corners] (0,0) -- (6,0) -- (6,4) -- (0,4) --  cycle;
    \begin{scope}
      \clip (0,0) -- (2,0)  .. controls (1.5,1.5) and (4.5,1.5)
      .. (4,4) -- (0,4) -- (0,0) -- cycle;
      \filldraw[fill=gray!40] (0,4) -- (0,2) .. controls (1.5,1.5)
      and (4.5,1.5) .. (4.5,0) -- (6,0) -- (6,4) -- (0,4) --  cycle;
    \end{scope}
    \draw [name path=line1] (2,0) .. controls (1.5,1.5) and (4.5,1.5) .. (4,4);
    \draw [name path=line2] (0,2)  .. controls (1.5,1.5) and (4.5,1.5) .. (4.5,0);
    \path [name intersections={of=line1 and line2}];
    \draw [rounded corners, use as bounding box, very thick] (0,0)
    -- (6,0) -- (6,4) -- (0,4) --  cycle;
    \node at (2,3) {$[w]$} ;
    \node at (1,1) {$[u]$} ;
    \node at (3,0.75) {$[v]$} ;
    \node at (4.75,2) {$[z]$} ;
  \end{tikzpicture}
  \caption{$W$ चा तुल्यतावर्गांमध्ये विभाग.}
  \ollabel{fig:partition}
\end{figure}

\end{document}
